% ============================================================================
% AION Autonomous Authority Acquisition, Systems Intelligence, and Delayed
% Outcome Apprenticeship
% Compile-ready section. Requires: amsmath, amssymb, booktabs, tabularx,
% array, url.
% ============================================================================

\section{Autonomous Authority Acquisition, Systems Intelligence, and Delayed-Outcome Apprenticeship}
\label{sec:aion-authority-systems-delayed-apprenticeship}

This development increment moved AION beyond one-off curriculum execution in
three linked directions.  First, it acquired reusable execution authorities
for real local toolchains and multi-project systems work.  Second, it applied
those authorities through a persistent executive that selected and completed
useful capability-acquisition objectives.  Third, it connected apprenticeship
and private repair to outcomes that AION did not generate, with credit withheld
until a later operational consequence became available.

The resulting control chain is
\begin{equation}
\begin{aligned}
\text{North-Star objective}
&\rightarrow \text{capability-gap and authority analysis} \\
&\rightarrow \text{private executor acquisition or useful-work selection} \\
&\rightarrow \text{precommitted plan, prediction, or artifact} \\
&\rightarrow \text{independently revealed consequence} \\
&\rightarrow \text{world-change versus internal-failure attribution} \\
&\rightarrow \text{private repair, abstention, or world-model revision} \\
&\rightarrow \text{later consequence confirmation} \\
&\rightarrow \text{executive receipt, curriculum evidence, and retention}.
\end{aligned}
\label{eq:aion-delayed-apprentice-chain}
\end{equation}

The essential authority rule is that proposal, execution, observation, and
promotion remain separate.  A tool adapter may propose evidence; it cannot
award a competency level.  A private repair may pass its own tests; it cannot
declare itself successful.  An artifact may contain a plausible conclusion; it
cannot verify that conclusion.  HexCore accepts each claim only through the
authority specified before the relevant outcome was observed.

\subsection{Reusable Progressive-Authority Acquisition}
\label{subsec:aion-progressive-authority-acquisition}

The progressive competency programme contained 44 leaf subjects, but many
could not progress because no suitable outcome authority existed.  Rather than
treat every subject as a bespoke benchmark, AION introduced an executor
registry and reusable authority-family acquisition layer.  For subject $s$ and
candidate adapter $e$, installation requires
\begin{equation}
\operatorname{Install}(e,s)=
\mathbf{1}\!\left[
  C_s\subseteq D_e
  \land X_e=1
  \land N_e=M_e
  \land U_e=0
  \land W_e=0
\right],
\label{eq:aion-adapter-install-gate}
\end{equation}
where $C_s$ is the declared subject subskill set, $D_e$ is the set actually
demonstrated by the executor, $X_e$ denotes successful execution, $N_e/M_e$ is
the rejected-counterexample ratio, $U_e$ is the number of unsafe variants
accepted, and $W_e$ is the number of live repository writes.

The executor registry records the adapter family, verified subskills, authority
artifact, acquisition status, and resume condition.  Missing adapters park the
affected subject and allow the curriculum to rotate; an unrelated executor is
never substituted merely to keep a progress counter moving.

\subsubsection{Progressive outcome-authority fabric}

The first factory generation converted five reusable outcome-authority families
into six verified subject adapters.  This was the bridge between the generic
registry and later local-toolchain acquisition.

\begin{table}[htbp]
\centering
\caption{Initial progressive outcome-authority fabric.}
\label{tab:aion-outcome-authority-fabric}
\begin{tabularx}{\textwidth}{@{}>{\raggedright\arraybackslash}p{4.0cm}>{\raggedright\arraybackslash}Xr@{}}
\toprule
Subject adapter & Outcome authority family & Faults \\
\midrule
Mathematics & exact proof and calculation & 8 \\
English natural-language cognition & document and delayed-question execution & 17 \\
Research and communication & document and delayed-question execution & 17 \\
Scientific method & empirical dataset or experiment & 15 \\
Data analysis and statistics & data and statistical outcome & 15 \\
Machine learning and AI & code and system execution & 10 \\
\bottomrule
\end{tabularx}
\end{table}

Each adapter had to execute a passing candidate, cover the declared subskills,
reject every mapped counterexample and six unsafe code classes, survive runtime
reconstruction, and perform no live write or objective mutation.  Unsupported
families remained explicit gaps.  The gate reported six verified adapters from
five distinct authority families, declared-subskill coverage, counterexample
falsification, unsupported-family abstention, restart retention, and zero unsafe
or objective mutations.  Verified executor coverage increased from $5/44$ to
$11/44$.

The promoted procedure is
\begin{center}
\small\path{procedure_autonomous_progressive_executor_factory_v1}.
\end{center}

\subsubsection{Governed cognitive self-improvement}

The same increment improved a live proposal component rather than only adding
subject adapters.  The retained mission-capability router was conservative: it
missed valid mixed mathematics--statistics--science objectives.  A permissive
alternative increased recall but routed irrelevant capabilities.  Four private
router policies therefore competed under a frozen champion, sealed mission
cohort, and protected single-domain retention set.

For mission $m$, expected capability set $Y_m$, and proposed set
$\widehat Y_m$, routing quality used set F1,
\begin{equation}
F_1(m)=
\frac{2\,|Y_m\cap\widehat Y_m|}
{|Y_m|+|\widehat Y_m|},
\label{eq:aion-router-set-f1}
\end{equation}
with exact-match, irrelevant-route, weakest-family, safety, and restart gates
evaluated separately.  The selected \texttt{discriminative\_transfer} policy
was the only challenger that improved mean routing while preserving the sealed
and protected floors.

\begin{table}[htbp]
\centering
\caption{Governed capability-router self-improvement.}
\label{tab:aion-router-self-improvement}
\begin{tabularx}{\textwidth}{@{}Xr@{}}
\toprule
Measure & Result \\
\midrule
Mean F1 improvement over retained control & $+11.11$ percentage points \\
Sealed mission exact routing & 100\% \\
Protected mission retention & 100\% \\
Malicious policies rejected & $5/5$ \\
Unsafe live writes / objective mutations & $0/0$ \\
Restart policy retention & passed \\
\bottomrule
\end{tabularx}
\end{table}

The promoted policy is content-hashed and loaded by the
\path{MissionCapabilityActionHarness}.  It remains proposal-only: it may suggest
which cognitive capabilities should handle a mission, but it cannot authorize
execution, award competence, rewrite the North Star, or bypass later outcome
verification.  The promoted procedure is
\begin{center}
\small\path{procedure_autonomous_cognitive_self_improvement_lab_v1}.
\end{center}

\subsubsection{Open local-toolchain acquisition}

AION inspected the local environment and discovered four usable execution
authorities: CPython, TypeScript plus Node, Rust, and SQLite.  It constructed a
private progressive examiner for each surface, tested a distinct faulty
implementation for every declared subskill, and installed only adapters that
passed the complete gate.  Go and Java were not installed because a valid local
runtime was unavailable.  This was treated as an explicit abstention rather
than simulated evidence.

\begin{table}[htbp]
\centering
\caption{Open toolchain acquisition result.}
\label{tab:aion-open-toolchain-acquisition}
\begin{tabularx}{\textwidth}{@{}lXr@{}}
\toprule
Authority & Verification role & Result \\
\midrule
CPython & isolated execution, behavioural checks, negative variants & installed \\
TypeScript + Node & strict compilation and runtime execution & installed \\
Rust & compilation and executable behavioural authority & installed \\
SQLite & schema, constraint, transaction, and query execution & installed \\
Go & required runtime not available & abstained \\
Java & required runtime not available & abstained \\
\midrule
\multicolumn{2}{l}{Deliberately faulty subskill implementations rejected} & $39/39$ \\
\multicolumn{2}{l}{Verified executor coverage} & $11/44\rightarrow15/44$ \\
\multicolumn{2}{l}{Unresolved catalogue gaps} & $33\rightarrow29$ \\
\multicolumn{2}{l}{Unsafe live writes / objective mutations} & $0/0$ \\
\bottomrule
\end{tabularx}
\end{table}

The promoted procedure is
\begin{center}
\small\path{procedure_open_toolchain_progressive_executor_acquisition_v1}.
\end{center}
It is an executor-acquisition result, not a claim that AION has mastered the complete
Python, TypeScript, Rust, or database ecosystems.

\subsection{Persistent Autonomous Useful-Work Selection}
\label{subsec:aion-useful-work-controller}

The executive layer was extended from method selection to iterative useful-work
ownership.  Candidate projects are ranked using explicit value, urgency, risk
reduction, learning value, cost, confidence, and dependency readiness.  The
implemented priority score is
\begin{equation}
P(i)=
\frac{
  \left(0.45V_i+0.25G_i+0.15R_i+0.15L_i\right)Q_i
}{\max(0.1,C_i)},
\label{eq:aion-useful-work-priority}
\end{equation}
where $V_i$ is value, $G_i$ urgency, $R_i$ risk reduction, $L_i$ learning
value, $Q_i$ confidence, and $C_i$ cost.  Ineligible dependency-blocked work
receives zero executable priority.

For every selected project, the controller must:
\begin{enumerate}
    \item preserve the constitutional parent objective;
    \item declare measurable success criteria before acting;
    \item compose an appropriate executive work system;
    \item execute only through a governed adapter;
    \item compare the resulting registry or operational state with the frozen
          baseline; and
    \item retain the next useful objective after restart.
\end{enumerate}

In its first generation the controller selected real toolchain acquisition,
rejected a low-value repeated suite, and parked a human-evaluation proposal
whose authority was unavailable.  In its second generation it remembered the
completed work, selected software/system-project authority acquisition, and
measured an 18.18 percentage-point executor-coverage gain.  The retained
procedure is \path{procedure_autonomous_useful_work_controller_v2}.

\subsection{Software and Systems Project Authority}
\label{subsec:aion-software-system-authority}

The next authority family replaced isolated language demonstrations with 17
distinct executable systems projects.  Each project contains a working
candidate, a semantically targeted fault mutation, mapped subskills, an
independent process outcome, and a security boundary.  The portfolio includes:
\begin{itemize}
    \item scoped capability tokens and supply-chain manifests;
    \item process supervision, filesystem permissions, and framed IPC;
    \item idempotent replicated logs and quorum decisions;
    \item progressive rollout, rollback, SLO monitoring, and capacity planning;
    \item container policy and architecture-boundary checks;
    \item incident containment, DNS caching, TLS contracts, and route failover;
    \item restartable operational handover.
\end{itemize}

The authority installed eight subject adapters:
\begin{center}
\small
Security Engineering; Operating Systems; Distributed Systems;
Cloud/DevOps/SRE; Architecture and Production Operations; Secure Software
Engineering; Computer Networking; and Integrated Engineering Capstone.
\end{center}

Evidence credit is intentionally narrower than the requested contract.  If
project $p$ declares tags $T_p$ and subject $s$ requires subskills $C_s$, the
credited set is
\begin{equation}
E_{p,s}=T_p\cap C_s,
\label{eq:aion-exact-subskill-credit}
\end{equation}
not the complete requested set.  This correction prevents a successful project
from silently awarding evidence for subskills it did not exercise.

\begin{table}[htbp]
\centering
\caption{Software/system authority promotion gate.}
\label{tab:aion-system-authority-gate}
\begin{tabularx}{\textwidth}{@{}Xr@{}}
\toprule
Measure & Result \\
\midrule
Verified progressive subject adapters & 8 \\
Distinct executable systems projects & 17 \\
Mapped project counterexamples rejected & $48/48$ \\
Declared subskill coverage & passed \\
Unsupported human authority & abstained \\
Executor coverage & $15/44\rightarrow23/44$ \\
Remaining executor gaps & $29\rightarrow21$ \\
Unsafe live writes / objective mutations & $0/0$ \\
\bottomrule
\end{tabularx}
\end{table}

The promoted authority is
\path{procedure_software_system_project_authority_v1}.  These projects are
engineered executable authorities; they enable progressive learning but do not
by themselves establish Advanced or Expert systems competence.

\subsection{Natural-Language Cross-System Capstone}
\label{subsec:aion-cross-system-capstone}

One broad objective required AION to design a secure distributed service that
survived interruption and deployment failure.  The mission-capability harness
selected exactly six interacting domains:
\begin{equation}
\mathcal{D}^{*}=\{
\text{distributed},\text{cloud/SRE},\text{architecture},
\text{operating systems},\text{networking},\text{security}
\}.
\end{equation}
The composed candidate integrated scoped cryptographic authorization,
idempotent replication, process-interruption recovery, TLS/network policy,
health-gated rollout, automatic rollback, and availability monitoring.

The evaluator separately tested the component transfers and the integrated
system.  It rejected four compound functional faults and six malicious
operational plans, including authority bypass and unsafe deployment behaviour.

\begin{table}[htbp]
\centering
\caption{Cross-system capstone result.}
\label{tab:aion-cross-system-capstone}
\begin{tabularx}{\textwidth}{@{}Xr@{}}
\toprule
Measure & Result \\
\midrule
Exact mission-to-domain route & passed \\
Source-disjoint component transfer & $6/6$ \\
Integrated execution & passed \\
Integrated counterexamples rejected & $4/4$ \\
Malicious plans rejected & $6/6$ \\
Unsafe live writes / objective mutations & $0/0$ \\
\bottomrule
\end{tabularx}
\end{table}

The promoted procedure is
\path{procedure_cross_system_intelligence_capstone_v1}.  This is a strong
internal multi-domain composition result, not proof of unrestricted autonomous
systems engineering.

\subsection{Consequence-Confirmed Natural Self-Repair}
\label{subsec:aion-consequence-confirmed-repair}

The private self-repair architecture was next connected to failures that had
actually occurred in AION's live progressive curriculum.  The repair watcher
reads two operational sources:
\begin{enumerate}
    \item the append-only curriculum executor attempt ledger; and
    \item the committed public-outcome campaign ledger.
\end{enumerate}
It does not inject a tutorial fault.  It detects a failure streak, hashes the
failure receipt, attributes the origin, and routes only internal faults into a
private executor trial.  A public revision changes the world model; it is not
evidence that AION itself is broken.

For outcome $o$, the implemented attribution order is
\begin{equation}
A(o)=
\begin{cases}
\text{persistence failure}, & \neg T(o),\\
\text{execution failure}, & T(o)\land\neg X(o),\\
\text{acquisition failure}, & \exists s:\neg R_s(o),\\
\text{external world change}, & \Delta_s(o)\neq0,\\
\text{stable observation}, & \text{otherwise},
\end{cases}
\label{eq:aion-outcome-attribution}
\end{equation}
where $T$ is transaction persistence, $X$ executable success, $R_s$ source
reachability, and $\Delta_s$ source revision change.

\paragraph{Two-stage repair authority.}
The private trial reconstructs the affected contract without writing
competency evidence.  It must execute successfully, reject its counterexample,
reject unsafe variants, and preserve the live champion.  Even then the repair
is only provisional.  It becomes consequence-confirmed only when the separate
curriculum service later produces a successful operational row for the same
contract and continues successfully on distinct contracts.

Repair identity excludes evidence that may accumulate later:
\begin{equation}
I_r=H(s,c,H(F),H(O_{\mathrm{confirm}})),
\label{eq:aion-stable-repair-identity}
\end{equation}
where $s$ is subject, $c$ contract, $F$ the immutable failure streak, and
$O_{\mathrm{confirm}}$ the first later confirmation.  Subsequent transfer
receipts append to $I_r$ without changing it, preventing double counting after
restart.

The operational ledger contained two closed episodes.  Algorithms and Data
Structures recorded 472 rejected attempts before its executable contract was
repaired and later confirmed.  Software Engineering failed three times, was
parked, and later passed after a verified executor repair.  These are
retrospective local operational incidents, not external production outages.

\begin{table}[htbp]
\centering
\caption{Consequence-confirmed self-repair result.}
\label{tab:aion-consequence-repair-result}
\begin{tabularx}{\textwidth}{@{}Xr@{}}
\toprule
Measure & Result \\
\midrule
Independently committed public outcomes processed & 14 \\
External world-change outcomes & 11 \\
External changes incorrectly routed to repair & 0 \\
Naturally occurring internal failure episodes & 2 \\
Private repairs confirmed by later outcomes & $2/2$ \\
Distinct internal domains & 2 \\
Later distinct-contract confirmations & 6 \\
Private repair validation rate & 100\% \\
Malicious repair strategies rejected & $6/6$ \\
Unsafe live writes / objective mutations & $0/0$ \\
\bottomrule
\end{tabularx}
\end{table}

The promoted procedure is
\begin{center}
\small\path{procedure_cross_domain_consequence_confirmed_self_repair_v1}.
\end{center}
Unresolved future failures are represented as pending repairs and cannot
contribute to the confirmed-repair count.

\subsection{Prospective Cross-Domain Public Outcomes}
\label{subsec:aion-prospective-public-outcomes}

The earlier long-duration campaign committed planned actions, but an action
plan is not a falsifiable forecast.  The prospective campaign therefore writes
a forecast commitment to an fsync-backed append-only ledger before invoking
any public authority.  The initial portfolio contains three independently
controlled authority families:
\begin{itemize}
    \item public source control: CPython and Node branch heads;
    \item public package registry: the published NumPy version; and
    \item public environmental data: Madrid temperature and wind from
          Open-Meteo.
\end{itemize}

For source-control and package observations, the one-step forecast is revision
stability,
\begin{equation}
\widehat r_{s,t+1}=r_{s,t}.
\end{equation}
For environmental observation $w_t=(\tau_t,v_t)$, the deliberately conservative
bounded-persistence forecast is
\begin{equation}
\tau_{t+1}\in[\tau_t-8,\tau_t+8],\qquad
v_{t+1}\leq\max(20,v_t+15).
\label{eq:aion-weather-persistence}
\end{equation}
The first observation establishes a baseline and cannot count as a correct
forecast.  If an authority is unreachable or omits a required field, the
forecast abstains from scoring rather than being treated as either correct or
incorrect.

Every observed row retains the commitment hash, public outcome, classification,
response, and final outcome hash.  Given $S$ scored forecasts, accuracy is
\begin{equation}
\operatorname{Acc}=\frac{1}{|S|}\sum_{j\in S}
\mathbf{1}[\widehat y_j=y_j].
\end{equation}
Promotion requires
\begin{equation}
G_{\mathrm{prospective}}=
\mathbf{1}[n\ge6]
\mathbf{1}[f\ge3]
\mathbf{1}[|S|\ge12]
\mathbf{1}[\operatorname{Acc}\ge0.75]
\mathbf{1}[t\ge1200\ \mathrm{s}]
\mathbf{1}[m=0]
\mathbf{1}[W=0],
\label{eq:aion-prospective-gate}
\end{equation}
where $n$ is cycle count, $f$ authority-family count, $t$ elapsed wall time,
$m$ external events misrouted to self-repair, and $W$ unsafe live writes.  The
wall-clock term prevents rapid replay from masquerading as continual
operation.

At the documented live snapshot, the campaign had eight prospective cycles,
28 scored forecasts, three authority families, 100\% accuracy, zero owner
interventions, and zero external events routed to self-repair.  It had run for
approximately 497 of the required 1200 seconds and was therefore correctly
reported as \texttt{COLLECTING\_PROSPECTIVE\_OUTCOMES}, not promoted.

\subsection{Owner-Valued Work Under Delayed Accountability}
\label{subsec:aion-owner-valued-delayed-work}

The public-outcome stream was then converted from passive monitoring into a
repeating useful-work campaign.  At generation $g$, AION derives an owner-facing
objective, fixes success criteria, composes an executive work system, produces
three artifacts, hashes them, and waits for public row $t+1$:
\begin{equation}
\begin{aligned}
O_t
&\rightarrow C_g
\rightarrow \{A_g^{R},A_g^{S},A_g^{D}\} \\
&\rightarrow H(C_g,A_g^{R},A_g^{S},A_g^{D})
\rightarrow O_{t+1}
\rightarrow \operatorname{Verify}_g.
\end{aligned}
\label{eq:aion-owner-work-generation}
\end{equation}
Here $A^R$ is a source-backed technical/release-risk brief, $A^S$ is an
executable public-outcome integrity guard, and $A^D$ is an environmental
monitoring decision.

\paragraph{Research lane.}
The brief records exact public source values, authority families, and a
provisional stability conclusion.  A later public repository/package row owns
the score.  With $q$ reachable claims,
\begin{equation}
\operatorname{ResearchPass}=
\mathbf{1}[q\ge2]
\mathbf{1}\!\left[
\frac{1}{q}\sum_{i=1}^{q}\mathbf{1}[\widehat r_i=r_i]\ge0.75
\right].
\end{equation}

\paragraph{Software lane.}
The generated integrity guard executes in a fresh isolated subprocess against
the later genuine row.  It recomputes the row hash, requires a precommitment
link, and requires at least three reachable authority families.  The verifier
then tampers with one outcome while preserving the claimed hash.  The software
lane passes only if
\begin{equation}
\operatorname{SoftwarePass}=
\mathbf{1}[\text{genuine later row accepted}]
\mathbf{1}[\text{tampered row rejected}].
\end{equation}

\paragraph{Data-decision lane.}
The monitoring decision fixes the same bounded environmental envelope before
the next weather row.  Missing authority causes explicit abstention.  A scored
decision passes only when the later measurement remains inside the committed
range.

Only after all three lanes pass is the generation recorded as a verified
executive method outcome.  Older evidence that predates executive-work-system
attachment remains valid provenance but does not receive retrospective
executive credit.

At the documented snapshot, three generations had produced nine artifacts.
Two generations had been confirmed by later public rows with 100\% success in
the research, software, and data lanes.  One qualifying executive receipt had
been recorded, and the third generation was waiting for its later consequence.
The campaign was therefore correctly reported as
\texttt{WAITING\_FOR\_LATER\_CONSEQUENCES}.  Promotion requires at least three
later-confirmed executive receipts and at least 75\% success in every lane.

\subsection{Runtime Integration and Persistent Supervision}
\label{subsec:aion-delayed-runtime-integration}

The long-lived real-outcome service now executes four connected operations in
order:
\begin{enumerate}
    \item ingest and arbitrate the existing long-duration public ledger;
    \item watch curriculum failures and close only consequence-confirmed
          private repairs;
    \item create the next prospective public forecast, durably commit it, and
          collect the independent outcome; and
    \item close the waiting owner-valued work generation and create its
          successor.
\end{enumerate}
The service repeats on a five-minute cadence under AION heartbeat supervision.
Runtime and dashboard reconstruction do not replay completed actions.  The
terminal dashboard exposes separate sections for consequence-confirmed repair,
prospective independent-outcome operation, and owner-valued delayed work.

\begin{table}[htbp]
\centering
\caption{Procedure and status summary at the documented snapshot.}
\label{tab:aion-delayed-procedure-status}
\begin{tabularx}{\textwidth}{@{}>{\raggedright\arraybackslash}p{5.1cm}>{\raggedright\arraybackslash}p{2.2cm}>{\raggedright\arraybackslash}X@{}}
\toprule
Procedure & Status & Principal evidence \\
\midrule
\path{procedure_autonomous_progressive_executor_factory_v1}
& promoted & 6 adapters across 5 reusable authority families \\
\path{procedure_autonomous_cognitive_self_improvement_lab_v1}
& promoted & $+11.11$ F1 points; sealed and protected routing 100\% \\
\path{procedure_open_toolchain_progressive_executor_acquisition_v1}
& promoted & 4 adapters; $39/39$ faults rejected \\
\path{procedure_autonomous_useful_work_controller_v2}
& promoted & sequential useful projects; persistent next objective \\
\path{procedure_software_system_project_authority_v1}
& promoted & 8 adapters; 17 projects; $48/48$ faults rejected \\
\path{procedure_cross_system_intelligence_capstone_v1}
& promoted & exact six-domain route; integrated execution and criticism \\
\path{procedure_cross_domain_consequence_confirmed_self_repair_v1}
& promoted & $2/2$ natural failures closed by later consequences \\
\path{procedure_prospective_cross_domain_outcome_operation_v1}
& collecting & cycle/forecast gates met; wall-clock gate still open \\
\path{procedure_owner_valued_delayed_work_v1}
& collecting & 2 generations confirmed; executive-receipt gate open \\
\bottomrule
\end{tabularx}
\end{table}

The complete focused integration stack passed 20 tests after the owner-valued
work integration.  These tests cover public-outcome precommitment, abstention,
later-row linkage, natural-failure identity, private repair, restart behaviour,
software tamper rejection, systems project authority, cross-system composition,
and dashboard reporting.

\subsection{Capability Meaning and Remaining Boundary}
\label{subsec:aion-autonomous-apprentice-boundary}

The evidence supports the following description:
\begin{quote}
AION is a governed, persistent, self-improving cognitive architecture with an
operational apprenticeship mechanism, expanding executable authorities, and a
live delayed-consequence learning loop.  It is becoming an autonomous general
apprentice; broad general apprenticeship has not yet been earned.
\end{quote}

The work does not support a claim of AGI.  Toolchain contracts, the 17 systems
projects, forecast forms, public-source portfolio, and the three owner-valued
work lanes remain engineered.  The public sources own their outcomes, but they
do not constitute unrestricted research, production CI, economic value, social
judgment, or physical-world breadth.

A capability-level autonomous-general-apprentice gate should require, at
minimum:
\begin{itemize}
    \item Advanced competence in at least six qualitatively different domains
          and Expert competence in at least two;
    \item at least 20 cross-domain projects with independent consequences;
    \item autonomous acquisition or invention of most missing executor families;
    \item several naturally occurring self-repairs verified only by later
          consequences;
    \item at least one retained cognitive-code improvement under a private
          tournament and later operational confirmation;
    \item retention over weeks, not merely process restart;
    \item declining owner intervention per new domain;
    \item repeated useful objectives generated and completed by the executive;
          and
    \item no serious backward forgetting, authority bypass, or unsafe live
          modification.
\end{itemize}

The next engineering boundary is open useful-objective and repair-contract
invention.  When a changing real portfolio exposes a capability for which no
retained executor, diagnostic, or repair vocabulary exists, AION must invent a
private authority contract, falsify that contract, produce owner-valued work,
and wait for a later independent consequence.  Further offline subject packs
are secondary unless they directly feed this repeated consequence loop.

\subsection{Principal Implementation Artifacts}
\label{subsec:aion-delayed-implementation-artifacts}

\begin{table}[htbp]
\centering
\caption{Principal implementation and evidence artifacts.}
\label{tab:aion-delayed-artifacts}
\begin{tabularx}{\textwidth}{@{}>{\raggedright\arraybackslash}p{3.5cm}>{\raggedright\arraybackslash}X@{}}
\toprule
Role & Repository path \\
\midrule
Authority-family factory & \path{backend/modules/hexcore/autonomous_progressive_executor_factory.py} \\
Cognitive improvement lab & \path{backend/modules/hexcore/autonomous_cognitive_self_improvement_lab.py} \\
Toolchain authority & \path{backend/modules/hexcore/open_toolchain_progressive_authority.py} \\
Useful-work controller & \path{backend/modules/hexcore/autonomous_useful_work_controller.py} \\
Systems authority & \path{backend/modules/hexcore/software_system_project_authority.py} \\
Cross-system capstone & \path{backend/modules/hexcore/cross_system_intelligence_capstone.py} \\
Consequence repair & \path{backend/modules/hexcore/cross_domain_consequence_repair.py} \\
Prospective campaign & \path{backend/modules/hexcore/prospective_cross_domain_outcome_campaign.py} \\
Owner-valued work & \path{backend/modules/hexcore/owner_valued_delayed_work_campaign.py} \\
Persistent service & \path{backend/AION/system/aion_real_outcome_learning_service.py} \\
Terminal dashboard & \path{backend/AION/system/aion_development_dashboard.py} \\
Consequence result & \path{results/hexcore_cross_domain_consequence_repair.json} \\
Prospective result & \path{results/hexcore_prospective_cross_domain_outcome.json} \\
Owner-work result & \path{results/hexcore_owner_valued_delayed_work.json} \\
\bottomrule
\end{tabularx}
\end{table}
